% SSSUP-FONTE-STORICA: fonti-locali/originali/2014-2020.zip
% Anno 2018, esercizio 1. Trasposizione dal sorgente LaTeX storico.
% Verificare editorialmente la trasposizione prima di considerarla definitiva.
\renewcommand{\Anno}{2018}
\renewcommand{\NumeroProblema}{1}

\renewcommand{\Testo}{
Lo stato di Santannaland utilizza come moneta nazionale il Piso e la Banca Centrale ha deciso di stampare banconote di soltanto due tagli. Si tenga conto che:
\begin{itemize}
\item tutti i beni hanno un prezzo intero, compreso tra 1 e 100 (estremi compresi);
\item i beni hanno prezzi  uniformemente ripartiti e la stessa probabilità di essere acquistati;
\item i pagamenti vengono fatti in contanti per la cifra esatta e senza ricevere resto;
\item ogni cittadino possiede banconote dei due tagli in gran quantità, sicuramente sufficienti a fare gli acquisti.
\end{itemize}

\begin{enumerate}
\item[(a)] Determinare la scelta dei due tagli di banconote che rende minimo in media il numero di banconote necessario per gli acquisti.
\item[(b)] Discutere la questione precedente nel caso generale in cui i prezzi dei beni siano compresi tra 1 ed $N$, con $N$ intero qualsiasi, lasciando uguale a due il numero di tagli di banconote da stampare.
\end{enumerate}
}

\renewcommand{\Risposta}{
% SSSUP-BOZZA-AI: Risposta; autore=ChatGPT; data=2026-11-12; stato=DA-VERIFICARE
% SSSUP-PROVENIENZA: ChatGPT (OpenAI), sessione ProblemiSSSUP del 2026-11-12; derivato dalla traccia canonica; revisione umana pendente.
Assumendo tagli interi positivi, (a) le due scelte ottimali sono
$1$ e $10$ Piso, oppure $1$ e $11$ Piso. In entrambi i casi
il numero medio minimo di banconote \`e $91/10=9{,}1$.

(b) Per $N\ge2$ un taglio deve essere $1$; il secondo \`e un intero
$d\in\{2,\ldots,N\}$ che minimizza
\[
 M_N(d)=\frac1N\sum_{s=1}^{N}
 \left(\left\lfloor\frac{s}{d}\right\rfloor+s-d\left\lfloor\frac{s}{d}\right\rfloor\right).
\]
La formula esplicita nella soluzione permette il confronto esatto
per ogni $N$, includendo eventuali ex aequo. Per $N=1$ basta il taglio $1$;
il secondo taglio, maggiore di $1$, non viene utilizzato.
}
\renewcommand{\Aiutino}{
% SSSUP-BOZZA-AI: Aiutino; autore=ChatGPT; data=2026-11-12; stato=DA-VERIFICARE
% SSSUP-PROVENIENZA: ChatGPT (OpenAI), sessione ProblemiSSSUP del 2026-11-12; derivato dalla traccia canonica; revisione umana pendente.
Quale taglio \`e obbligatorio per pagare un bene da un Piso?
Fissato l'altro taglio $d$, scrivere ciascun prezzo come $s=ad+b$,
con $0\le b<d$. Per calcolare la media, raggruppare i prezzi
in blocchi di $d$ valori consecutivi. Non dare per scontato
che il taglio ottimale sia unico.
}
\renewcommand{\SoluzioneUfficiale}{}
\renewcommand{\SoluzioneRedazionale}{
% SSSUP-BOZZA-AI: SoluzioneRedazionale; autore=ChatGPT; data=2026-11-12; stato=DA-VERIFICARE
% SSSUP-PROVENIENZA: ChatGPT (OpenAI), sessione ProblemiSSSUP del 2026-11-12; derivato dalla traccia canonica; revisione umana pendente.
\textbf{Convenzione sui tagli.}
Interpretiamo i tagli come numeri interi positivi di Piso e i prezzi
$1,\ldots,N$ come equiprobabili. L'integralit\`a dei tagli, usuale
in questa lettura del problema, non \`e esplicitata nel Testo.

Per poter pagare esattamente un Piso, uno dei due tagli deve essere $1$.
Indichiamo l'altro con $d\ge2$. Per pagare $s$, conviene usare il maggior
numero possibile di banconote da $d$: sostituire una di queste con $d$
banconote da $1$ aumenta infatti il numero di banconote di $d-1$.
Scrivendo $s=ad+b$, con $0\le b<d$, il numero minimo \`e quindi $a+b$.

\textbf{Formula per la somma.}
Poniamo $N=qd+r$, con $0\le r<d$, e indichiamo con $S_N(d)$ la somma
dei numeri minimi di banconote per tutti i prezzi da $1$ a $N$.
Poich\'e
\[
 \sum_{s=1}^{N}\left\lfloor\frac{s}{d}\right\rfloor
 =\frac{dq(q-1)}2+q(r+1),
\]
si ottiene
\[
 \boxed{
 S_N(d)=\frac{N(N+1)}2
 -(d-1)\left(\frac{dq(q-1)}2+q(r+1)\right),
 \qquad M_N(d)=\frac{S_N(d)}N.
 }
\]
Nella somma dei quozienti ci sono $d$ occorrenze di ciascun valore
$0,\ldots,q-1$ e $r+1$ occorrenze di $q$, se si include anche il prezzo
$0$, che non altera il risultato.

\textbf{(a) Il caso $N=100$.}
Basta considerare $2\le d\le100$: un taglio maggiore di $100$ non viene
mai usato e d\`a media $50{,}5$.
Se $d\le5$, il numero di banconote per pagare $s$ \`e almeno $s/d$,
e quindi
\[
 S_{100}(d)\ge\frac{5050}{d}\ge1010.
\]
Se $d\ge21$, ogni possibile resto modulo $d$ compare al massimo cinque
volte fra i prezzi $1,\ldots,100$. La somma dei cento resti \`e pertanto
almeno $5(0+1+\cdots+19)=950$; a maggior ragione
$S_{100}(d)\ge950$.
Restano soltanto quindici valori, per i quali la formula precedente d\`a
\[
\begin{array}{c|rrrrrrrr}
 d&6&7&8&9&10&11&12&13\\\hline
 S_{100}(d)&1050&976&934&914&910&910&914&934
\end{array}
\]
\[
\begin{array}{c|rrrrrrr}
 d&14&15&16&17&18&19&20\\\hline
 S_{100}(d)&955&976&1000&1050&1055&1090&1155.
\end{array}
\]
Il minimo \`e $910$, raggiunto soltanto per $d=10$ e $d=11$.
La media minima \`e dunque $910/100=9{,}1$.

\textbf{(b) Il caso generale.}
Per $N\ge2$ non serve considerare $d>N$, perch\'e il taglio non verrebbe
usato, mentre $d=2$ migliora strettamente tale scelta.
La formula di $S_N(d)$ fornisce quindi un procedimento finito ed esatto:
per ciascun $d=2,\ldots,N$ si calcolano quoziente $q$, resto $r$ e
$S_N(d)$, e si conservano tutti i valori di $d$ che realizzano il minimo.
Questa \`e una caratterizzazione esatta per ogni $N$, non un'approssimazione
basata sull'arrotondamento di una radice quadrata.
L'ottimo pu\`o non essere unico, come gi\`a mostra il caso $N=100$.
Per $N=1$ occorre soltanto una banconota da $1$ per l'unico acquisto.
}
\renewcommand{\SoluzioniAlternative}{}
\renewcommand{\Osservazioni}{}
\renewcommand{\Commenti}{}
